\PassOptionsToPackage{unicode}{hyperref}
\PassOptionsToPackage{hyphens}{url}
\PassOptionsToPackage{dvipsnames,svgnames*,x11names*}{xcolor}
\documentclass[
  11pt,
]{article}
\usepackage{amsmath,amssymb}
\usepackage{lmodern}
\usepackage{iftex}
\ifPDFTeX
  \usepackage[T1]{fontenc}
  \usepackage[utf8]{inputenc}
  \usepackage{textcomp}
\else
  \usepackage{unicode-math}
  \defaultfontfeatures{Scale=MatchLowercase}
  \defaultfontfeatures[\rmfamily]{Ligatures=TeX,Scale=1}
\fi
\IfFileExists{upquote.sty}{\usepackage{upquote}}{}
\IfFileExists{microtype.sty}{
  \usepackage[]{microtype}
  \UseMicrotypeSet[protrusion]{basicmath}
}{}
\makeatletter
\@ifundefined{KOMAClassName}{
  \IfFileExists{parskip.sty}{
    \usepackage{parskip}
  }{
    \setlength{\parindent}{0pt}
    \setlength{\parskip}{6pt plus 2pt minus 1pt}}
}{
  \KOMAoptions{parskip=half}}
\makeatother
\usepackage{xcolor}
\IfFileExists{xurl.sty}{\usepackage{xurl}}{}
\IfFileExists{bookmark.sty}{\usepackage{bookmark}}{\usepackage{hyperref}}
\hypersetup{
  colorlinks=true,
  linkcolor={blue},
  filecolor={Maroon},
  citecolor={Blue},
  urlcolor={blue},
  pdfcreator={LaTeX via pandoc}}
\urlstyle{same}
\usepackage[margin=2cm]{geometry}
\usepackage{color}
\usepackage{fancyvrb}
\newcommand{\VerbBar}{|}
\newcommand{\VERB}{\Verb[commandchars=\\\{\}]}
\DefineVerbatimEnvironment{Highlighting}{Verbatim}{commandchars=\\\{\}}
\newenvironment{Shaded}{}{}
\newcommand{\AlertTok}[1]{\textcolor[rgb]{1.00,0.00,0.00}{\textbf{#1}}}
\newcommand{\AnnotationTok}[1]{\textcolor[rgb]{0.38,0.63,0.69}{\textbf{\textit{#1}}}}
\newcommand{\AttributeTok}[1]{\textcolor[rgb]{0.49,0.56,0.16}{#1}}
\newcommand{\BaseNTok}[1]{\textcolor[rgb]{0.25,0.63,0.44}{#1}}
\newcommand{\BuiltInTok}[1]{#1}
\newcommand{\CharTok}[1]{\textcolor[rgb]{0.25,0.44,0.63}{#1}}
\newcommand{\CommentTok}[1]{\textcolor[rgb]{0.38,0.63,0.69}{\textit{#1}}}
\newcommand{\CommentVarTok}[1]{\textcolor[rgb]{0.38,0.63,0.69}{\textbf{\textit{#1}}}}
\newcommand{\ConstantTok}[1]{\textcolor[rgb]{0.53,0.00,0.00}{#1}}
\newcommand{\ControlFlowTok}[1]{\textcolor[rgb]{0.00,0.44,0.13}{\textbf{#1}}}
\newcommand{\DataTypeTok}[1]{\textcolor[rgb]{0.56,0.13,0.00}{#1}}
\newcommand{\DecValTok}[1]{\textcolor[rgb]{0.25,0.63,0.44}{#1}}
\newcommand{\DocumentationTok}[1]{\textcolor[rgb]{0.73,0.13,0.13}{\textit{#1}}}
\newcommand{\ErrorTok}[1]{\textcolor[rgb]{1.00,0.00,0.00}{\textbf{#1}}}
\newcommand{\ExtensionTok}[1]{#1}
\newcommand{\FloatTok}[1]{\textcolor[rgb]{0.25,0.63,0.44}{#1}}
\newcommand{\FunctionTok}[1]{\textcolor[rgb]{0.02,0.16,0.49}{#1}}
\newcommand{\ImportTok}[1]{#1}
\newcommand{\InformationTok}[1]{\textcolor[rgb]{0.38,0.63,0.69}{\textbf{\textit{#1}}}}
\newcommand{\KeywordTok}[1]{\textcolor[rgb]{0.00,0.44,0.13}{\textbf{#1}}}
\newcommand{\NormalTok}[1]{#1}
\newcommand{\OperatorTok}[1]{\textcolor[rgb]{0.40,0.40,0.40}{#1}}
\newcommand{\OtherTok}[1]{\textcolor[rgb]{0.00,0.44,0.13}{#1}}
\newcommand{\PreprocessorTok}[1]{\textcolor[rgb]{0.74,0.48,0.00}{#1}}
\newcommand{\RegionMarkerTok}[1]{#1}
\newcommand{\SpecialCharTok}[1]{\textcolor[rgb]{0.25,0.44,0.63}{#1}}
\newcommand{\SpecialStringTok}[1]{\textcolor[rgb]{0.73,0.40,0.53}{#1}}
\newcommand{\StringTok}[1]{\textcolor[rgb]{0.25,0.44,0.63}{#1}}
\newcommand{\VariableTok}[1]{\textcolor[rgb]{0.10,0.09,0.49}{#1}}
\newcommand{\VerbatimStringTok}[1]{\textcolor[rgb]{0.25,0.44,0.63}{#1}}
\newcommand{\WarningTok}[1]{\textcolor[rgb]{0.38,0.63,0.69}{\textbf{\textit{#1}}}}
\usepackage{longtable,booktabs,array}
\usepackage{calc}
\usepackage{etoolbox}
\makeatletter
\patchcmd\longtable{\par}{\if@noskipsec\mbox{}\fi\par}{}{}
\makeatother
\IfFileExists{footnotehyper.sty}{\usepackage{footnotehyper}}{\usepackage{footnote}}
\makesavenoteenv{longtable}
\setlength{\emergencystretch}{3em}
\providecommand{\tightlist}{
  \setlength{\itemsep}{0pt}\setlength{\parskip}{0pt}}
\setcounter{secnumdepth}{5}
\usepackage{fvextra}
\DefineVerbatimEnvironment{Highlighting}{Verbatim}{breaklines,breakanywhere,fontsize=\footnotesize,commandchars=\\\{\}}
\RecustomVerbatimEnvironment{verbatim}{Verbatim}{breaklines,breakanywhere,fontsize=\footnotesize}
\usepackage{microtype}
\usepackage{booktabs}
\usepackage{longtable}
\usepackage{array}
\usepackage{enumitem}
\setlist[itemize]{itemsep=2pt,parsep=2pt}
\setlist[enumerate]{itemsep=2pt,parsep=2pt}
\usepackage{etoolbox}
\AtBeginEnvironment{longtable}{\footnotesize}
\AtBeginEnvironment{tabular}{\footnotesize}
\renewcommand{\arraystretch}{1.2}
\setlength{\tabcolsep}{4pt}
\usepackage{fancyhdr}
\usepackage{titlesec}
\titleformat{\section}{\Large\bfseries}{\thesection}{0.6em}{}
\titleformat{\subsection}{\large\bfseries}{\thesubsection}{0.6em}{}
\titlespacing*{\section}{0pt}{1.4em}{0.6em}
\titlespacing*{\subsection}{0pt}{1.0em}{0.4em}
\setlength{\parindent}{0pt}
\setlength{\parskip}{0.5em plus 0.1em minus 0.1em}
\pagestyle{fancy}
\fancyhf{}
\fancyhead[L]{\nouppercase{\leftmark}}
\fancyhead[R]{\thepage}
\fancyfoot[C]{overlap-calculator case studies}
\renewcommand{\headrulewidth}{0.4pt}
\renewcommand{\footrulewidth}{0pt}
\setlength{\headheight}{14pt}
\sloppy
\ifLuaTeX
  \usepackage{selnolig}
\fi

\title{Case Study 05 — Default Light Source Selection}
\usepackage{etoolbox}
\makeatletter
\providecommand{\subtitle}[1]{
  \apptocmd{\@title}{\par {\large #1 \par}}{}{}
}
\makeatother
\subtitle{Subset of bundled built-in sources}
\author{}
\date{}

\begin{document}
\maketitle

{
\hypersetup{linkcolor=}
\setcounter{tocdepth}{2}
\tableofcontents
}
\hypertarget{case-study-05--default-light-source-selection}{
\section{Case Study 05 --- Default Light Source
Selection}\label{case-study-05--default-light-source-selection}}

Pick a \textbf{subset} of the bundled built-in light sources rather than
running against the full default set. Here we run against \texttt{LEDB1}
and \texttt{LEDB4} --- the former is \emph{not} in the default set, so
it is never analysed unless you opt into it explicitly.

The user-selectable knob is:

\begin{itemize}
\tightlist
\item
  CLI: \texttt{-\/-default-light-sources\ NAME{[},NAME...{]}}
\item
  API: form field \texttt{default\_light\_sources=NAME{[},NAME...{]}}
\end{itemize}

Allowed names: \texttt{AM15G}, \texttt{LEDB1}, \texttt{LEDB2},
\texttt{LEDB3}, \texttt{LEDB4}, \texttt{CIEFL10}.

\begin{center}\rule{0.5\linewidth}{0.5pt}\end{center}

\hypertarget{what-you-will-produce}{
\subsection{What you will produce}\label{what-you-will-produce}}

Everything is written under one \texttt{output/} directory inside this
case study folder:

\begin{longtable}[]{@{}lll@{}}
\toprule
Folder & By & Contents \\
\midrule
\endhead
\texttt{output/cli/tables/} & CLI & Result tables (CSV / JSON / XLSX) \\
\texttt{output/cli/plots/} & CLI & TIFF plots per sample, light source,
and (sample, light) pair \\
\texttt{output/api/} & API & Same \texttt{tables/} and \texttt{plots/}
tree, extracted from the ZIP \\
\texttt{output/api/analysis\_outputs.zip} & API & Raw ZIP response \\
\bottomrule
\end{longtable}

Light sources used: \textbf{LED-B1 + LED-B4}.

\begin{center}\rule{0.5\linewidth}{0.5pt}\end{center}

\hypertarget{1-prerequisites}{
\subsection{1. Prerequisites}\label{1-prerequisites}}

If you have never run any case study before, follow
\href{https://github.com/pinarsyt/overlap-calculator/blob/main/case_studies/01_theoretical_only/README.md\#1-prerequisites}{Case
Study 01 §1 Prerequisites} once to install the package, the
\texttt{requests} library, and either the Python or Docker option for
the HTTP API. Every case study uses the same environment.

\begin{center}\rule{0.5\linewidth}{0.5pt}\end{center}

\hypertarget{2-inputs}{
\subsection{2. Inputs}\label{2-inputs}}

The case-study \texttt{input.json} lists 2 Gaussian TD-DFT outputs. The
manifest uses relative paths that point at the bundled data under
\href{https://github.com/pinarsyt/overlap-calculator/blob/main/input/files}{\texttt{../../input/files/}},
so no files are copied or duplicated. To swap in your own data, replace
\texttt{input.json} with a manifest of the same shape; you do not need
to touch \texttt{submit.py}.

\hypertarget{3-run-with-the-cli}{
\subsection{3. Run with the CLI}\label{3-run-with-the-cli}}

Open a terminal at the repository root, then run a single command:

\begin{Shaded}
\begin{Highlighting}[]
\ExtensionTok{overlap{-}calculator}\NormalTok{ analyze }\AttributeTok{{-}{-}input}\NormalTok{ case\_studies/05\_default\_light\_source\_selection/input.json }\AttributeTok{{-}{-}out}\NormalTok{ case\_studies/05\_default\_light\_source\_selection/output/cli }\AttributeTok{{-}{-}default{-}light{-}sources}\NormalTok{ LEDB1,LEDB4 }\AttributeTok{{-}{-}plot{-}dpi}\NormalTok{ 400}
\end{Highlighting}
\end{Shaded}

After the run,
\texttt{case\_studies/05\_default\_light\_source\_selection/output/cli/}
contains a \texttt{tables/} directory (and a \texttt{plots/} directory
unless plots were disabled). The \texttt{-\/-default-light-sources} flag
pins the bundled sources to LED-B1 + LED-B4; remove it to compare
against the full default set (\texttt{AM15G,LEDB4,LEDB2,LEDB3,CIEFL10}).

\begin{center}\rule{0.5\linewidth}{0.5pt}\end{center}

\hypertarget{4-run-with-the-http-api}{
\subsection{4. Run with the HTTP API}\label{4-run-with-the-http-api}}

Start the API as documented in
\href{https://github.com/pinarsyt/overlap-calculator/blob/main/case_studies/01_theoretical_only/README.md\#1-prerequisites}{Case
Study 01 §1} and verify with
\texttt{curl\ http://localhost:8000/health} (or \texttt{:8080} if you
started the API via \texttt{docker\ compose}). Then submit the same
uploads in one curl command:

\begin{Shaded}
\begin{Highlighting}[]
\ExtensionTok{curl} \AttributeTok{{-}X}\NormalTok{ POST http://localhost:8000/analyze }\AttributeTok{{-}F} \StringTok{"files=@input/files/slurm{-}5473089.out"} \AttributeTok{{-}F} \StringTok{"files=@input/files/slurm{-}5473178.out"} \AttributeTok{{-}F} \StringTok{"default\_light\_sources=LEDB1,LEDB4"} \AttributeTok{{-}F} \StringTok{"plot\_dpi=400"} \AttributeTok{{-}{-}output}\NormalTok{ case\_studies/05\_default\_light\_source\_selection/output/api/analysis\_outputs.zip}
\end{Highlighting}
\end{Shaded}

Each \texttt{-F\ "files=@..."} is one upload. The form fields mirror the
CLI flags. The response is a ZIP archive saved as
\texttt{output/api/analysis\_outputs.zip}.

\begin{center}\rule{0.5\linewidth}{0.5pt}\end{center}

\hypertarget{5-drive-both-interfaces-from-one-python-script}{
\subsection{5. Drive both interfaces from one Python
script}\label{5-drive-both-interfaces-from-one-python-script}}

Instead of typing the commands above, run
\href{https://github.com/pinarsyt/overlap-calculator/blob/main/case_studies/05_default_light_source_selection/submit.py}{\texttt{submit.py}}
from the repository root:

\begin{Shaded}
\begin{Highlighting}[]
\ExtensionTok{python}\NormalTok{ case\_studies/05\_default\_light\_source\_selection/submit.py}
\end{Highlighting}
\end{Shaded}

The script reads \texttt{input.json}, runs the CLI step (writing to
\texttt{output/cli/}), then POSTs the same uploads to \texttt{/analyze}
and unpacks the response into \texttt{output/api/}. If the API is not
reachable, it prints a hint and exits cleanly without failing the CLI
step.

\begin{center}\rule{0.5\linewidth}{0.5pt}\end{center}

\hypertarget{6-expected-output-layout}{
\subsection{6. Expected output layout}\label{6-expected-output-layout}}

\begin{verbatim}
case_studies/05_default_light_source_selection/output/
+-- cli/
|   +-- tables/results.{csv,json,xlsx}        (2 samples × 2 light sources = 4 rows)
|   +-- tables/results_timings.{csv,json,xlsx}
|   +-- tables/descriptor_summary.{csv,xlsx}
|   +-- tables/ranking_by_light_source__<metric>.{csv,json,xlsx}   (gaussian/lorentzian x absorbed_fraction/shape_overlap)
|   +-- tables/ranking_by_sample__<metric>.{csv,json,xlsx}          (same four metric variants)
|   +-- plots/                                                      (also plots/ranking/by_light_source/<metric>/ and plots/ranking/by_sample/<metric>/)
+-- api/
    +-- analysis_outputs.zip
    +-- tables/...
    +-- plots/...
\end{verbatim}

API sample IDs are auto-generated from each upload's \texttt{\%chk}
header, while the CLI uses the \texttt{sample\_id} you set in
\texttt{input.json}. This is the only cosmetic difference between the
two trees; the tables align row by row otherwise.

\begin{center}\rule{0.5\linewidth}{0.5pt}\end{center}

\hypertarget{7-related-case-studies}{
\subsection{7. Related case studies}\label{7-related-case-studies}}

\begin{itemize}
\tightlist
\item
  \href{https://github.com/pinarsyt/overlap-calculator/blob/main/case_studies/04_custom_light_source/README.md}{Case
  Study 04 --- Custom Light Source}
\item
  \href{https://github.com/pinarsyt/overlap-calculator/blob/main/case_studies/01_theoretical_only/README.md}{Case
  Study 01 --- Theoretical Only}
\end{itemize}

\end{document}
